\documentclass[11pt]{article}
% =====================================================================
%  Gemeinsame Vorlage der Arbeitsblätter – Künstliche Intelligenz (Q13, gA LB 4 / eA LB 2)
%  Kopf, Fuß und Antwortplatz wie informatik/13/rechner/arbeitsblaetter/ab-vorlage.tex
%  (Linien 8 mm, Karo 5 mm aus ganzen Kästchen, Lücken seit 2026-10 x 1,5).
%  Einbinden in jedem Blatt direkt nach \documentclass: % =====================================================================
%  Gemeinsame Vorlage der Arbeitsblätter – Künstliche Intelligenz (Q13, gA LB 4 / eA LB 2)
%  Kopf, Fuß und Antwortplatz wie informatik/13/rechner/arbeitsblaetter/ab-vorlage.tex
%  (Linien 8 mm, Karo 5 mm aus ganzen Kästchen, Lücken seit 2026-10 x 1,5).
%  Einbinden in jedem Blatt direkt nach \documentclass: % =====================================================================
%  Gemeinsame Vorlage der Arbeitsblätter – Künstliche Intelligenz (Q13, gA LB 4 / eA LB 2)
%  Kopf, Fuß und Antwortplatz wie informatik/13/rechner/arbeitsblaetter/ab-vorlage.tex
%  (Linien 8 mm, Karo 5 mm aus ganzen Kästchen, Lücken seit 2026-10 x 1,5).
%  Einbinden in jedem Blatt direkt nach \documentclass: \input{ab-vorlage}
%  Übersetzen mit XeLaTeX (oder Tectonic):
%
%  Blatt:   xelatex ab-3-k-means.tex
%  Lösung:  xelatex ab-3-k-means-lsg.tex
%           (Datei enthält nur: \def\mitloesung{}\input{ab-3-k-means})
%  Lösungseinträge belegen nie mehr Platz als die Lücke im Blatt.
%
%  Antwortplatz:
%   \luecke[l] / \lsg[l]{Text}      Lücke (Standard 6,75 cm, Fachbegriff 5,25 cm, Zahl 2,25 cm)
%   \restzeile / \lsgrest{Text}     Lücke bis Zeilenende      \lueckentext  Absatz mit Lücken
%   \linien{n} / \lsglinien{n}{T}   Schreiblinien 8 mm        \linienfeld{n} / \lsgfeld{n}{Code}
%   \schreibzeile                   Tabellenzelle 8 mm        \lsgtext{…}  nur in der Lösung
%   \kariert{n} / \lsgkariert{n}{TikZ}, \karobox{b}{h} / \lsgkarobox{b}{h}{TikZ}
%   \kreuz{0|1}                     Ankreuzfeld (1 = in der Lösung angekreuzt)
%  KI:
%   java / javaluecke               Listings (Lücken mit (*\lsg[3cm]{…}*))
%   \nr{1}                          Schrittnummer ① wie auf der Website und in der IDE
%   \niveauEA                       Kennzeichen „eA“ für Aufgaben des erhöhten Niveaus
%   Farben: adelie, zuegel, esels (Pinguinarten), klnull, kleins (Klassen 0/1), cA … cF (Cluster)
% =====================================================================
\usepackage[a4paper,left=16mm,right=16mm,top=23mm,bottom=13mm,headheight=12mm,headsep=3mm,footskip=6mm]{geometry}
\usepackage{iftex}
\ifPDFTeX
  \usepackage[T1]{fontenc}
  \usepackage[utf8]{inputenc}
  \usepackage{helvet}
  \usepackage[scaled=0.86]{beramono}
  \renewcommand{\familydefault}{\sfdefault}
\else
  \usepackage{fontspec}
  \setmainfont{texgyreheros}[Extension=.otf,UprightFont=*-regular,BoldFont=*-bold,ItalicFont=*-italic,BoldItalicFont=*-bolditalic]
  \setmonofont{DejaVuSansMono}[Extension=.ttf,UprightFont=*,BoldFont=*-Bold,ItalicFont=*-Oblique,BoldItalicFont=*-BoldOblique,Scale=0.86]
\fi
\usepackage[ngerman,shorthands=off]{babel}
\usepackage{amssymb}
\usepackage{xcolor,graphicx,tabularx,array,enumitem,fancyhdr,tikz,multicol,calc,etoolbox,pifont,xspace,listings}
\usepackage[most]{tcolorbox}
\usetikzlibrary{arrows.meta,positioning,calc,shapes.geometric,decorations.pathreplacing}
\usepackage[hidelinks]{hyperref}
\setlength{\parindent}{0pt}
\setlength{\parskip}{3pt}
\setlist{nosep,leftmargin=*}

% ---------- Lösungsschalter ----------
\newif\ifloesung
\ifdefined\mitloesung\loesungtrue\fi
\newcommand{\lsgfarbe}{\color{red!75!black}}

% ---------- Kopf- und Fußzeile ----------
\newcommand{\lehrkraft}{Hau}
\newcommand{\themenbereich}{Künstliche Intelligenz}
\newcommand{\abnummer}{}
\pagestyle{fancy}\fancyhf{}
\renewcommand{\headrulewidth}{0.4pt}
\fancyhead[L]{\small Klasse 13\\Datum:}
\fancyhead[C]{\small Themenbereich:\\\themenbereich}
\fancyhead[R]{\small \lehrkraft\ifloesung\\{\lsgfarbe\bfseries Lösung}\fi}
\fancyfoot[C]{\scriptsize edu-mrh.de\,·\,Informatik 13\,·\,Blatt \abnummer\ifloesung\ (Lösung)\fi\,·\,Seite \thepage}
\newcommand{\abtitel}[2]{\renewcommand{\abnummer}{#1}\begin{center}\Large\bfseries #1\quad\underline{#2}\end{center}\vspace{-2mm}}

% ---------- Boxen ----------
\newenvironment{aufgabe}[2][]{\begin{tcolorbox}[enhanced,breakable,colback=white,colframe=black,boxrule=0.7pt,arc=3mm,left=2mm,right=2mm,top=1mm,bottom=1.2mm,before skip=2mm,after skip=2mm]\textbf{Aufgabe #2}\ifstrempty{#1}{}{ \textit{#1}}\par\vspace{0.3mm}}{\end{tcolorbox}}
\newtcolorbox{merke}[1][Merke]{enhanced,breakable,colback=green!5,colframe=green!40!black,boxrule=0.8pt,arc=2mm,left=2mm,right=2mm,top=1mm,bottom=1.2mm,before skip=2mm,after skip=2mm,title={#1},fonttitle=\bfseries,coltitle=white,toptitle=0.3mm,bottomtitle=0.3mm}
\newtcolorbox{hinweis}{enhanced,breakable,colback=yellow!12,colframe=orange!70!black,boxrule=0.6pt,arc=2mm,left=2mm,right=2mm,top=1mm,bottom=1mm,before skip=2mm,after skip=2mm}
\newcommand{\web}[1]{{\small\color{gray}\textit{Website: #1}}}
\newcommand{\code}[1]{\texttt{#1}}

% ---------- Antwortplatz (wie informatik/10/java, Lücken x 1,5) ----------
\newlength{\zeilenabstand}\setlength{\zeilenabstand}{8mm}
\newlength{\karo}\setlength{\karo}{5mm}
\tikzset{karo/.style={draw=gray!50,line width=0.3pt},schreiblinie/.style={draw=black!65,line width=0.35pt}}
\newlength{\lsgbreite}
\newcommand{\lsg}[2][6.75cm]{\leavevmode\ifloesung\settowidth{\lsgbreite}{\,#2\,}%
  \ifdim\lsgbreite>#1\relax\rlap{\raisebox{0.8pt}{\resizebox{#1}{!}{\lsgfarbe\,#2\,}}}%
  \else\rlap{\raisebox{0.8pt}{\lsgfarbe\,#2}}\fi\fi\rule[-0.6ex]{#1}{0.4pt}\xspace}
\newcommand{\luecke}[1][6.75cm]{\lsg[#1]{}}
\newcommand{\lsgrest}[1]{\leavevmode\unskip\hspace{1.5mm}\ifloesung\rlap{\raisebox{0.8pt}{\lsgfarbe\,#1}}\fi\leaders\hbox{\rule[-0.6ex]{0.5pt}{0.4pt}}\hfill\null\par}
\newcommand{\restzeile}{\lsgrest{}}
\newcommand{\lsgtext}[1]{\ifloesung\leavevmode{\lsgfarbe#1}\fi}
\newcommand{\lsgoder}[2]{\ifloesung\leavevmode#1\else#2\fi}
\newenvironment{lueckentext}{\par\linespread{1.5}\selectfont\raggedright}{\par}
\newcommand{\schreibzeile}{\rule[-2.8mm]{0pt}{8mm}}
% Tabellenzelle für n Zeilen Handschrift (Lösungstext darin kleiner setzen: \lsgklein{…})
\newcommand{\schreibzeilen}[1]{\rule[\dimexpr-#1\zeilenabstand+5.2mm\relax]{0pt}{\dimexpr#1\zeilenabstand\relax}}
\newcommand{\lsgklein}[1]{\lsgtext{\footnotesize\raggedright #1}}
\newcommand{\nr}[1]{\ding{\numexpr171+#1\relax}}
\newcommand{\lsglinien}[2]{\par\vspace{0.5mm}\noindent\begin{tikzpicture}
  \path (0,0) (\linewidth-0.4pt,-#1\zeilenabstand-1.5mm);
  \foreach \i in {1,...,#1}{\draw[schreiblinie] (0.2pt,-\i\zeilenabstand) -- (\linewidth-0.2pt,-\i\zeilenabstand);}
  \ifloesung\node[overlay,anchor=north west,inner sep=0pt] at (1mm,-\zeilenabstand+1pt+5mm) {\parbox[t]{\linewidth-2mm}{\lsgfarbe\raggedright\setlength{\baselineskip}{\zeilenabstand}\rule{0pt}{5mm}#2}};\fi
  \end{tikzpicture}\par}
\newcommand{\linien}[1]{\lsglinien{#1}{}}
\newcommand{\lsgfeld}[2]{\noindent\begin{tikzpicture}
  \draw[line width=0.4pt] (0.2pt,0) rectangle (\linewidth-0.2pt,-#1\zeilenabstand-3mm);
  \foreach \i in {1,...,#1}{\draw[schreiblinie] (2mm,-\i\zeilenabstand) -- (\linewidth-2mm,-\i\zeilenabstand);}
  \ifloesung\node[overlay,anchor=north west,inner sep=0pt] at (2.5mm,-\zeilenabstand+1pt+5mm) {\parbox[t]{\linewidth-5mm}{\lsgfarbe\ttfamily\small\setlength{\baselineskip}{\zeilenabstand}\rule{0pt}{5mm}#2}};\fi
  \end{tikzpicture}}
\newcommand{\linienfeld}[1]{\lsgfeld{#1}{}}
\newcommand{\karoraster}[3]{\begin{tikzpicture}\draw[karo] (0,0) grid[step=\karo] (#1\karo,#2\karo);
  \ifloesung\begin{scope}[overlay,shift={(0,#2\karo)},red!75!black,text=red!75!black,line width=0.8pt]#3\end{scope}\fi\end{tikzpicture}}
\newcommand{\lsgkariert}[2]{\par\vspace{1.5mm}\noindent\pgfmathtruncatemacro{\karoX}{floor(\linewidth/\karo)}%
  \makebox[\linewidth]{\karoraster{\karoX}{#1}{#2}}\par\vspace{1mm}}
\newcommand{\kariert}[1]{\lsgkariert{#1}{}}
\newcommand{\lsgkarobox}[3]{\pgfmathtruncatemacro{\karoX}{round((#1)/\karo)}\pgfmathtruncatemacro{\karoY}{round((#2)/\karo)}\karoraster{\karoX}{\karoY}{#3}}
\newcommand{\karobox}[2]{\lsgkarobox{#1}{#2}{}}
% Ankreuzfeld: \kreuz{1} = angekreuzt (nur in der Lösung)
\newcommand{\kreuz}[1]{\leavevmode\ifloesung\ifnum#1=1 {\lsgfarbe$\boxtimes$}\else$\square$\fi\else$\square$\fi}

% ---------- Listings: Java und Python ----------
\lstset{basicstyle=\ttfamily\small,keywordstyle=\bfseries,commentstyle=\color{gray!80!black},stringstyle=\color{black},showstringspaces=false,columns=flexible,keepspaces=true,tabsize=3,frame=single,framerule=0.4pt,rulecolor=\color{black},xleftmargin=2mm,framexleftmargin=1mm,aboveskip=1.5mm,belowskip=1.5mm,breaklines=true,escapeinside={(*}{*)},
 literate={ä}{{\"a}}1 {ö}{{\"o}}1 {ü}{{\"u}}1 {Ä}{{\"A}}1 {Ö}{{\"O}}1 {Ü}{{\"U}}1 {ß}{{\ss}}1 {–}{{--}}1 {…}{{\dots}}1 {≤}{{$\le$}}1 {≥}{{$\ge$}}1 {≠}{{$\ne$}}1 {·}{{$\cdot$}}1 {①}{{\nr{1}}}1 {②}{{\nr{2}}}1 {③}{{\nr{3}}}1 {④}{{\nr{4}}}1}
\lstnewenvironment{java}[1][]{\lstset{language=Java,#1}}{}
\lstnewenvironment{javaluecke}[1][]{\lstset{language=Java,basicstyle=\linespread{1.5}\small\ttfamily,#1}}{}
\lstnewenvironment{python}[1][]{\lstset{language=Python,#1}}{}
\lstnewenvironment{pythonluecke}[1][]{\lstset{language=Python,basicstyle=\linespread{1.5}\small\ttfamily,#1}}{}
\lstnewenvironment{prolog}[1][]{\lstset{morecomment=[l]{\%},#1}}{}
\lstnewenvironment{prologluecke}[1][]{\lstset{morecomment=[l]{\%},basicstyle=\linespread{1.5}\small\ttfamily,#1}}{}

\newcommand{\niveauEA}{\tikz[baseline=(t.base)]\node[fill=orange!20,draw=orange!70!black,rounded corners=1.5mm,inner sep=1pt,font=\scriptsize\bfseries,text=orange!50!black](t){\,eA\,};}

% ---------- TikZ-Stile ----------
\tikzset{
  neuron/.style={circle,draw=black,line width=0.7pt,minimum size=8mm,inner sep=0pt,font=\small},
  eingabe/.style={neuron,fill=gray!12},
  kante/.style={-{Stealth[length=2mm]},line width=0.6pt},
  gew/.style={font=\scriptsize,fill=white,inner sep=0.6pt},
  pt/.style={circle,inner sep=0pt,minimum size=2.2mm,draw=white,line width=0.3pt},
}
% ---------- Farben wie auf der Website ----------
\definecolor{adelie}{HTML}{EA580C}\definecolor{zuegel}{HTML}{9333EA}\definecolor{esels}{HTML}{0E7490}
\definecolor{klnull}{HTML}{2563EB}\definecolor{kleins}{HTML}{DC2626}
\definecolor{cA}{HTML}{2563EB}\definecolor{cB}{HTML}{DC2626}\definecolor{cC}{HTML}{16A34A}\definecolor{cD}{HTML}{D97706}\definecolor{cE}{HTML}{9333EA}\definecolor{cF}{HTML}{0891B2}
\setlength{\unitlength}{1mm}

%  Übersetzen mit XeLaTeX (oder Tectonic):
%
%  Blatt:   xelatex ab-3-k-means.tex
%  Lösung:  xelatex ab-3-k-means-lsg.tex
%           (Datei enthält nur: \def\mitloesung{}\documentclass[11pt]{article}
\input{ab-vorlage}
\begin{document}
\abtitel{3}{k-Means: Ordnung ohne Lehrer}

\begin{aufgabe}[k-Means von Hand, $k=3$]{1}
\begin{minipage}[t]{0.6\linewidth}\vspace{0pt}
\begin{tikzpicture}[x=5mm,y=5mm]
  \draw[karo] (0,0) grid (19,19);
  \draw[-{Stealth[length=2mm]},line width=0.6pt] (0,0) -- (19.6,0) node[below left,font=\footnotesize]{$x$};
  \draw[-{Stealth[length=2mm]},line width=0.6pt] (0,0) -- (0,19.6) node[below left,font=\footnotesize]{$y$};
  \foreach \v in {1,...,9}{\node[font=\tiny,below] at (2*\v,0) {\v}; \node[font=\tiny,left] at (0,2*\v) {\v};}
  \foreach \n/\x/\y in {A/1/2,B/2/1,C/2/3,D/3/2,E/7/8,F/8/7,G/8/9,H/9/8,I/6/2,J/8/1,K/9/3,L/5/6}{
    \fill (2*\x,2*\y) circle (2.4pt); \node[font=\scriptsize,above right,inner sep=1pt] at (2*\x,2*\y) {\n};}
  \ifloesung
    \begin{scope}[red!75!black,line width=0.9pt]
      \foreach \x/\y/\t in {2/3/1,9/8/2,8/1/3}{\draw (2*\x-0.5,2*\y-0.5) -- (2*\x+0.5,2*\y+0.5) (2*\x-0.5,2*\y+0.5) -- (2*\x+0.5,2*\y-0.5);}
      \draw[dashed,-{Stealth}] (4,6) -- (5.2,5.6); \draw[dashed,-{Stealth}] (5.2,5.6) -- (4,4);
      \draw[dashed,-{Stealth}] (18,16) -- (16,16); \draw[dashed,-{Stealth}] (16,16) -- (14.8,15.2);
      \draw[dashed,-{Stealth}] (16,2) -- (15.33,4);
      \foreach \x/\y in {4/4,14.8/15.2,15.33/4}{\draw[line width=1.3pt] (\x-0.55,\y-0.55) -- (\x+0.55,\y+0.55) (\x-0.55,\y+0.55) -- (\x+0.55,\y-0.55);}
    \end{scope}
  \fi
\end{tikzpicture}
\end{minipage}\hfill
\begin{minipage}[t]{0.38\linewidth}\vspace{0pt}
Startzentren: Z1 = C(2\,\textbar\,3), Z2 = H(9\,\textbar\,8), Z3 = J(8\,\textbar\,1).\par\vspace{1mm}
\nr{1}~Jeder Punkt geht zum nächsten Zentrum.\\
\nr{2}~Jedes Zentrum wandert in die Mitte seiner Punkte.\\
\nr{3}~Wiederholen, bis sich nichts mehr ändert.\par\vspace{1.5mm}
Zeichne die Zentren als Kreuz ein und ihre Wege als Pfeile.\par\vspace{1.5mm}
Nur bei L ist es knapp. Berechne im 2.~Durchgang die Abstände von L zu Z1 und Z2:
\begin{lueckentext}
$d(L,Z1) = $ \lsg[4.2cm]{$\sqrt{2{,}4^2+3{,}2^2}=4$}\\
$d(L,Z2) = $ \lsg[4.2cm]{$\sqrt{3^2+2^2}\approx 3{,}61$}\\
L gehört dann zu \lsg[2.25cm]{Z2}.
\end{lueckentext}
\end{minipage}\par\vspace{2mm}
\renewcommand{\arraystretch}{1.1}
\begin{tabularx}{\linewidth}{|c|X|X|X|p{4.7cm}|}\hline
\textbf{Durchgang} & \textbf{Punkte bei Z1} & \textbf{Punkte bei Z2} & \textbf{Punkte bei Z3} & \textbf{neue Zentren (\nr{2})} \\\hline
1\schreibzeilen{2} & \lsgklein{A, B, C, D, L} & \lsgklein{E, F, G, H} & \lsgklein{I, J, K} & \lsgklein{Z1(2,6\,|\,2,8), Z2(8\,|\,8),\\ Z3(7,67\,|\,2)} \\\hline
2\schreibzeilen{2} & \lsgklein{A, B, C, D} & \lsgklein{E, F, G, H, L} & \lsgklein{I, J, K} & \lsgklein{Z1(2\,|\,2), Z2(7,4\,|\,7,6),\\ Z3(7,67\,|\,2)} \\\hline
3\schreibzeilen{1} & \lsgklein{A, B, C, D} & \lsgklein{E, F, G, H, L} & \lsgklein{I, J, K} & \lsgklein{keine Änderung: fertig} \\\hline
\end{tabularx}
\end{aufgabe}

\begin{aufgabe}[Der Algorithmus (Website: Festhalten)]{2}
\begin{lueckentext}
k-Means ist ein Verfahren des \lsg[5.25cm]{unüberwachten} Lernens: Die Punkte haben keine \lsg[5.25cm]{Etiketten (Klassen)}. Vorgegeben wird nur die Zahl \lsg[2.25cm]{$k$} der Cluster. \nr{1}~Jeder Punkt wird dem Zentrum mit dem kleinsten \lsg[4.5cm]{Abstand} zugeordnet. \nr{2}~Jedes Zentrum wird auf den \lsg[5.25cm]{Mittelwert} seiner Punkte gesetzt. \nr{3}~Beides wird wiederholt, bis sich \lsgrest{beim Zuordnen nichts mehr ändert.}
Die Streuung (Summe der quadrierten Abstände zum eigenen Zentrum) wird dabei nie \lsg[4.5cm]{größer}.
\end{lueckentext}
\end{aufgabe}

\newpage
\begin{aufgabe}[Selbst bauen (Online-IDE, TODO 3b und 3d)]{3}
Ergänze die Lücken in den Kernzeilen von \code{zuordnen()} und \code{lauf()}.
\begin{javaluecke}
for (int i = 0; i < punkte.length; i++) {                    // (*\nr{1}*) zuordnen()
   int bester = 0;  double kleinster = abstand(punkte[i].getWerte(), zentren[0]);
   for (int c = 1; c < k; c++) {
      double d = abstand(punkte[i].getWerte(), (*\lsg[2.8cm]{zentren[c]}*));
      if ((*\lsg[3.6cm]{d < kleinster}*)) { kleinster = d;  bester = c; }
   }
   if (punkte[i].getCluster() != bester) { (*\lsg[5.2cm]{punkte[i].setCluster(bester);}*) geaendert = true; }
}
while ((*\lsg[2.8cm]{zuordnen()}*)) { (*\lsg[4.4cm]{zentrenVerschieben();}*) durchgaenge++; }     // (*\nr{3}*) lauf()
\end{javaluecke}
\end{aufgabe}

\begin{aufgabe}[Startpunkte (Online-IDE)]{4}
\renewcommand{\arraystretch}{1.1}
\begin{tabularx}{\linewidth}{|p{2.8cm}|p{1.9cm}|p{1.9cm}|X|}\hline
\textbf{Start} & \textbf{Durchg.} & \textbf{Streuung} & \textbf{Was fällt auf?} \\\hline
\{0, 200, 300\}\schreibzeilen{2} & \lsgklein{3} & \lsgklein{13,55} & \lsgklein{Esels fast perfekt in einem Cluster; Adélie und Zügel nur teilweise getrennt (11 Adélie bei den Zügeln).} \\\hline
\{150, 160, 170\}\schreibzeilen{2} & \lsgklein{5} & \lsgklein{18,41} & \lsgklein{Esels auf zwei Cluster geteilt, Adélie und Zügel zusammen: schlechtere Lösung (lokales Minimum).} \\\hline
\end{tabularx}\par\vspace{1mm}
Wie kann man vermeiden, dass k-Means in einer schlechten Lösung hängen bleibt?
\lsglinien{1}{Mehrmals mit verschiedenen Starts rechnen, das Ergebnis mit der kleinsten Streuung nehmen.}
\end{aufgabe}

\begin{aufgabe}[Wie viele Cluster? Der Ellbogen (Online-IDE, TODO 5)]{5}
\begin{minipage}[t]{0.5\linewidth}\vspace{0pt}
\begin{tikzpicture}[x=5mm,y=5mm]
  \draw[karo] (0,0) grid (14,10);
  \draw[-{Stealth[length=2mm]},line width=0.6pt] (0,0) -- (14.6,0) node[below left,font=\footnotesize]{$k$};
  \draw[-{Stealth[length=2mm]},line width=0.6pt] (0,0) -- (0,10.6) node[right,font=\footnotesize]{Streuung};
  \foreach \v in {1,...,6}{\node[font=\tiny,below] at (2*\v,0) {\v};}
  \foreach \v in {10,20,...,50}{\node[font=\tiny,left] at (0,\v/5.5) {\v};}
  \ifloesung
    \draw[red!75!black,line width=0.9pt] (2,9.39) -- (4,3.81) -- (6,2.46) -- (8,2.01) -- (10,1.59) -- (12,1.36);
    \foreach \x/\y in {2/9.39,4/3.81,6/2.46,8/2.01,10/1.59,12/1.36}{\fill[red!75!black] (\x,\y) circle (2.2pt);}
  \fi
\end{tikzpicture}
\end{minipage}\hfill
\begin{minipage}[t]{0.47\linewidth}\vspace{0pt}
Trage die Streuung für $k = 1$ bis 6 ein und verbinde die Punkte.
\begin{lueckentext}
Der Knick liegt bei $k = $ \lsg[2.25cm]{3}.
Begründung: \lsgrest{Danach bringt jedes}
\lsgrest{weitere Cluster nur noch wenig.}
\end{lueckentext}
Warum ist das $k$ mit der kleinsten Streuung nicht das beste?
\lsglinien{2}{Sie sinkt mit jedem Cluster; bei k = 342 wäre sie 0 (jeder Pinguin allein): nutzlos.}
\end{minipage}
\end{aufgabe}

\begin{aufgabe}[Grenzen (Website)]{6}
Bei den „zwei Monden“ und beim „Ring um Kern“ findet k-Means nicht die Gruppen, die ein Mensch sieht. Erkläre mit der Lage der Grenze zwischen zwei Zentren.
\lsglinien{2}{Die Grenze zwischen zwei Zentren ist eine Gerade (Mittelsenkrechte). Gebogene oder ineinanderliegende Gruppen trennt sie nicht; k-Means passt zu runden Gruppen.}
\end{aufgabe}
\end{document}
)
%  Lösungseinträge belegen nie mehr Platz als die Lücke im Blatt.
%
%  Antwortplatz:
%   \luecke[l] / \lsg[l]{Text}      Lücke (Standard 6,75 cm, Fachbegriff 5,25 cm, Zahl 2,25 cm)
%   \restzeile / \lsgrest{Text}     Lücke bis Zeilenende      \lueckentext  Absatz mit Lücken
%   \linien{n} / \lsglinien{n}{T}   Schreiblinien 8 mm        \linienfeld{n} / \lsgfeld{n}{Code}
%   \schreibzeile                   Tabellenzelle 8 mm        \lsgtext{…}  nur in der Lösung
%   \kariert{n} / \lsgkariert{n}{TikZ}, \karobox{b}{h} / \lsgkarobox{b}{h}{TikZ}
%   \kreuz{0|1}                     Ankreuzfeld (1 = in der Lösung angekreuzt)
%  KI:
%   java / javaluecke               Listings (Lücken mit (*\lsg[3cm]{…}*))
%   \nr{1}                          Schrittnummer ① wie auf der Website und in der IDE
%   \niveauEA                       Kennzeichen „eA“ für Aufgaben des erhöhten Niveaus
%   Farben: adelie, zuegel, esels (Pinguinarten), klnull, kleins (Klassen 0/1), cA … cF (Cluster)
% =====================================================================
\usepackage[a4paper,left=16mm,right=16mm,top=23mm,bottom=13mm,headheight=12mm,headsep=3mm,footskip=6mm]{geometry}
\usepackage{iftex}
\ifPDFTeX
  \usepackage[T1]{fontenc}
  \usepackage[utf8]{inputenc}
  \usepackage{helvet}
  \usepackage[scaled=0.86]{beramono}
  \renewcommand{\familydefault}{\sfdefault}
\else
  \usepackage{fontspec}
  \setmainfont{texgyreheros}[Extension=.otf,UprightFont=*-regular,BoldFont=*-bold,ItalicFont=*-italic,BoldItalicFont=*-bolditalic]
  \setmonofont{DejaVuSansMono}[Extension=.ttf,UprightFont=*,BoldFont=*-Bold,ItalicFont=*-Oblique,BoldItalicFont=*-BoldOblique,Scale=0.86]
\fi
\usepackage[ngerman,shorthands=off]{babel}
\usepackage{amssymb}
\usepackage{xcolor,graphicx,tabularx,array,enumitem,fancyhdr,tikz,multicol,calc,etoolbox,pifont,xspace,listings}
\usepackage[most]{tcolorbox}
\usetikzlibrary{arrows.meta,positioning,calc,shapes.geometric,decorations.pathreplacing}
\usepackage[hidelinks]{hyperref}
\setlength{\parindent}{0pt}
\setlength{\parskip}{3pt}
\setlist{nosep,leftmargin=*}

% ---------- Lösungsschalter ----------
\newif\ifloesung
\ifdefined\mitloesung\loesungtrue\fi
\newcommand{\lsgfarbe}{\color{red!75!black}}

% ---------- Kopf- und Fußzeile ----------
\newcommand{\lehrkraft}{Hau}
\newcommand{\themenbereich}{Künstliche Intelligenz}
\newcommand{\abnummer}{}
\pagestyle{fancy}\fancyhf{}
\renewcommand{\headrulewidth}{0.4pt}
\fancyhead[L]{\small Klasse 13\\Datum:}
\fancyhead[C]{\small Themenbereich:\\\themenbereich}
\fancyhead[R]{\small \lehrkraft\ifloesung\\{\lsgfarbe\bfseries Lösung}\fi}
\fancyfoot[C]{\scriptsize edu-mrh.de\,·\,Informatik 13\,·\,Blatt \abnummer\ifloesung\ (Lösung)\fi\,·\,Seite \thepage}
\newcommand{\abtitel}[2]{\renewcommand{\abnummer}{#1}\begin{center}\Large\bfseries #1\quad\underline{#2}\end{center}\vspace{-2mm}}

% ---------- Boxen ----------
\newenvironment{aufgabe}[2][]{\begin{tcolorbox}[enhanced,breakable,colback=white,colframe=black,boxrule=0.7pt,arc=3mm,left=2mm,right=2mm,top=1mm,bottom=1.2mm,before skip=2mm,after skip=2mm]\textbf{Aufgabe #2}\ifstrempty{#1}{}{ \textit{#1}}\par\vspace{0.3mm}}{\end{tcolorbox}}
\newtcolorbox{merke}[1][Merke]{enhanced,breakable,colback=green!5,colframe=green!40!black,boxrule=0.8pt,arc=2mm,left=2mm,right=2mm,top=1mm,bottom=1.2mm,before skip=2mm,after skip=2mm,title={#1},fonttitle=\bfseries,coltitle=white,toptitle=0.3mm,bottomtitle=0.3mm}
\newtcolorbox{hinweis}{enhanced,breakable,colback=yellow!12,colframe=orange!70!black,boxrule=0.6pt,arc=2mm,left=2mm,right=2mm,top=1mm,bottom=1mm,before skip=2mm,after skip=2mm}
\newcommand{\web}[1]{{\small\color{gray}\textit{Website: #1}}}
\newcommand{\code}[1]{\texttt{#1}}

% ---------- Antwortplatz (wie informatik/10/java, Lücken x 1,5) ----------
\newlength{\zeilenabstand}\setlength{\zeilenabstand}{8mm}
\newlength{\karo}\setlength{\karo}{5mm}
\tikzset{karo/.style={draw=gray!50,line width=0.3pt},schreiblinie/.style={draw=black!65,line width=0.35pt}}
\newlength{\lsgbreite}
\newcommand{\lsg}[2][6.75cm]{\leavevmode\ifloesung\settowidth{\lsgbreite}{\,#2\,}%
  \ifdim\lsgbreite>#1\relax\rlap{\raisebox{0.8pt}{\resizebox{#1}{!}{\lsgfarbe\,#2\,}}}%
  \else\rlap{\raisebox{0.8pt}{\lsgfarbe\,#2}}\fi\fi\rule[-0.6ex]{#1}{0.4pt}\xspace}
\newcommand{\luecke}[1][6.75cm]{\lsg[#1]{}}
\newcommand{\lsgrest}[1]{\leavevmode\unskip\hspace{1.5mm}\ifloesung\rlap{\raisebox{0.8pt}{\lsgfarbe\,#1}}\fi\leaders\hbox{\rule[-0.6ex]{0.5pt}{0.4pt}}\hfill\null\par}
\newcommand{\restzeile}{\lsgrest{}}
\newcommand{\lsgtext}[1]{\ifloesung\leavevmode{\lsgfarbe#1}\fi}
\newcommand{\lsgoder}[2]{\ifloesung\leavevmode#1\else#2\fi}
\newenvironment{lueckentext}{\par\linespread{1.5}\selectfont\raggedright}{\par}
\newcommand{\schreibzeile}{\rule[-2.8mm]{0pt}{8mm}}
% Tabellenzelle für n Zeilen Handschrift (Lösungstext darin kleiner setzen: \lsgklein{…})
\newcommand{\schreibzeilen}[1]{\rule[\dimexpr-#1\zeilenabstand+5.2mm\relax]{0pt}{\dimexpr#1\zeilenabstand\relax}}
\newcommand{\lsgklein}[1]{\lsgtext{\footnotesize\raggedright #1}}
\newcommand{\nr}[1]{\ding{\numexpr171+#1\relax}}
\newcommand{\lsglinien}[2]{\par\vspace{0.5mm}\noindent\begin{tikzpicture}
  \path (0,0) (\linewidth-0.4pt,-#1\zeilenabstand-1.5mm);
  \foreach \i in {1,...,#1}{\draw[schreiblinie] (0.2pt,-\i\zeilenabstand) -- (\linewidth-0.2pt,-\i\zeilenabstand);}
  \ifloesung\node[overlay,anchor=north west,inner sep=0pt] at (1mm,-\zeilenabstand+1pt+5mm) {\parbox[t]{\linewidth-2mm}{\lsgfarbe\raggedright\setlength{\baselineskip}{\zeilenabstand}\rule{0pt}{5mm}#2}};\fi
  \end{tikzpicture}\par}
\newcommand{\linien}[1]{\lsglinien{#1}{}}
\newcommand{\lsgfeld}[2]{\noindent\begin{tikzpicture}
  \draw[line width=0.4pt] (0.2pt,0) rectangle (\linewidth-0.2pt,-#1\zeilenabstand-3mm);
  \foreach \i in {1,...,#1}{\draw[schreiblinie] (2mm,-\i\zeilenabstand) -- (\linewidth-2mm,-\i\zeilenabstand);}
  \ifloesung\node[overlay,anchor=north west,inner sep=0pt] at (2.5mm,-\zeilenabstand+1pt+5mm) {\parbox[t]{\linewidth-5mm}{\lsgfarbe\ttfamily\small\setlength{\baselineskip}{\zeilenabstand}\rule{0pt}{5mm}#2}};\fi
  \end{tikzpicture}}
\newcommand{\linienfeld}[1]{\lsgfeld{#1}{}}
\newcommand{\karoraster}[3]{\begin{tikzpicture}\draw[karo] (0,0) grid[step=\karo] (#1\karo,#2\karo);
  \ifloesung\begin{scope}[overlay,shift={(0,#2\karo)},red!75!black,text=red!75!black,line width=0.8pt]#3\end{scope}\fi\end{tikzpicture}}
\newcommand{\lsgkariert}[2]{\par\vspace{1.5mm}\noindent\pgfmathtruncatemacro{\karoX}{floor(\linewidth/\karo)}%
  \makebox[\linewidth]{\karoraster{\karoX}{#1}{#2}}\par\vspace{1mm}}
\newcommand{\kariert}[1]{\lsgkariert{#1}{}}
\newcommand{\lsgkarobox}[3]{\pgfmathtruncatemacro{\karoX}{round((#1)/\karo)}\pgfmathtruncatemacro{\karoY}{round((#2)/\karo)}\karoraster{\karoX}{\karoY}{#3}}
\newcommand{\karobox}[2]{\lsgkarobox{#1}{#2}{}}
% Ankreuzfeld: \kreuz{1} = angekreuzt (nur in der Lösung)
\newcommand{\kreuz}[1]{\leavevmode\ifloesung\ifnum#1=1 {\lsgfarbe$\boxtimes$}\else$\square$\fi\else$\square$\fi}

% ---------- Listings: Java und Python ----------
\lstset{basicstyle=\ttfamily\small,keywordstyle=\bfseries,commentstyle=\color{gray!80!black},stringstyle=\color{black},showstringspaces=false,columns=flexible,keepspaces=true,tabsize=3,frame=single,framerule=0.4pt,rulecolor=\color{black},xleftmargin=2mm,framexleftmargin=1mm,aboveskip=1.5mm,belowskip=1.5mm,breaklines=true,escapeinside={(*}{*)},
 literate={ä}{{\"a}}1 {ö}{{\"o}}1 {ü}{{\"u}}1 {Ä}{{\"A}}1 {Ö}{{\"O}}1 {Ü}{{\"U}}1 {ß}{{\ss}}1 {–}{{--}}1 {…}{{\dots}}1 {≤}{{$\le$}}1 {≥}{{$\ge$}}1 {≠}{{$\ne$}}1 {·}{{$\cdot$}}1 {①}{{\nr{1}}}1 {②}{{\nr{2}}}1 {③}{{\nr{3}}}1 {④}{{\nr{4}}}1}
\lstnewenvironment{java}[1][]{\lstset{language=Java,#1}}{}
\lstnewenvironment{javaluecke}[1][]{\lstset{language=Java,basicstyle=\linespread{1.5}\small\ttfamily,#1}}{}
\lstnewenvironment{python}[1][]{\lstset{language=Python,#1}}{}
\lstnewenvironment{pythonluecke}[1][]{\lstset{language=Python,basicstyle=\linespread{1.5}\small\ttfamily,#1}}{}
\lstnewenvironment{prolog}[1][]{\lstset{morecomment=[l]{\%},#1}}{}
\lstnewenvironment{prologluecke}[1][]{\lstset{morecomment=[l]{\%},basicstyle=\linespread{1.5}\small\ttfamily,#1}}{}

\newcommand{\niveauEA}{\tikz[baseline=(t.base)]\node[fill=orange!20,draw=orange!70!black,rounded corners=1.5mm,inner sep=1pt,font=\scriptsize\bfseries,text=orange!50!black](t){\,eA\,};}

% ---------- TikZ-Stile ----------
\tikzset{
  neuron/.style={circle,draw=black,line width=0.7pt,minimum size=8mm,inner sep=0pt,font=\small},
  eingabe/.style={neuron,fill=gray!12},
  kante/.style={-{Stealth[length=2mm]},line width=0.6pt},
  gew/.style={font=\scriptsize,fill=white,inner sep=0.6pt},
  pt/.style={circle,inner sep=0pt,minimum size=2.2mm,draw=white,line width=0.3pt},
}
% ---------- Farben wie auf der Website ----------
\definecolor{adelie}{HTML}{EA580C}\definecolor{zuegel}{HTML}{9333EA}\definecolor{esels}{HTML}{0E7490}
\definecolor{klnull}{HTML}{2563EB}\definecolor{kleins}{HTML}{DC2626}
\definecolor{cA}{HTML}{2563EB}\definecolor{cB}{HTML}{DC2626}\definecolor{cC}{HTML}{16A34A}\definecolor{cD}{HTML}{D97706}\definecolor{cE}{HTML}{9333EA}\definecolor{cF}{HTML}{0891B2}
\setlength{\unitlength}{1mm}

%  Übersetzen mit XeLaTeX (oder Tectonic):
%
%  Blatt:   xelatex ab-3-k-means.tex
%  Lösung:  xelatex ab-3-k-means-lsg.tex
%           (Datei enthält nur: \def\mitloesung{}\documentclass[11pt]{article}
% =====================================================================
%  Gemeinsame Vorlage der Arbeitsblätter – Künstliche Intelligenz (Q13, gA LB 4 / eA LB 2)
%  Kopf, Fuß und Antwortplatz wie informatik/13/rechner/arbeitsblaetter/ab-vorlage.tex
%  (Linien 8 mm, Karo 5 mm aus ganzen Kästchen, Lücken seit 2026-10 x 1,5).
%  Einbinden in jedem Blatt direkt nach \documentclass: \input{ab-vorlage}
%  Übersetzen mit XeLaTeX (oder Tectonic):
%
%  Blatt:   xelatex ab-3-k-means.tex
%  Lösung:  xelatex ab-3-k-means-lsg.tex
%           (Datei enthält nur: \def\mitloesung{}\input{ab-3-k-means})
%  Lösungseinträge belegen nie mehr Platz als die Lücke im Blatt.
%
%  Antwortplatz:
%   \luecke[l] / \lsg[l]{Text}      Lücke (Standard 6,75 cm, Fachbegriff 5,25 cm, Zahl 2,25 cm)
%   \restzeile / \lsgrest{Text}     Lücke bis Zeilenende      \lueckentext  Absatz mit Lücken
%   \linien{n} / \lsglinien{n}{T}   Schreiblinien 8 mm        \linienfeld{n} / \lsgfeld{n}{Code}
%   \schreibzeile                   Tabellenzelle 8 mm        \lsgtext{…}  nur in der Lösung
%   \kariert{n} / \lsgkariert{n}{TikZ}, \karobox{b}{h} / \lsgkarobox{b}{h}{TikZ}
%   \kreuz{0|1}                     Ankreuzfeld (1 = in der Lösung angekreuzt)
%  KI:
%   java / javaluecke               Listings (Lücken mit (*\lsg[3cm]{…}*))
%   \nr{1}                          Schrittnummer ① wie auf der Website und in der IDE
%   \niveauEA                       Kennzeichen „eA“ für Aufgaben des erhöhten Niveaus
%   Farben: adelie, zuegel, esels (Pinguinarten), klnull, kleins (Klassen 0/1), cA … cF (Cluster)
% =====================================================================
\usepackage[a4paper,left=16mm,right=16mm,top=23mm,bottom=13mm,headheight=12mm,headsep=3mm,footskip=6mm]{geometry}
\usepackage{iftex}
\ifPDFTeX
  \usepackage[T1]{fontenc}
  \usepackage[utf8]{inputenc}
  \usepackage{helvet}
  \usepackage[scaled=0.86]{beramono}
  \renewcommand{\familydefault}{\sfdefault}
\else
  \usepackage{fontspec}
  \setmainfont{texgyreheros}[Extension=.otf,UprightFont=*-regular,BoldFont=*-bold,ItalicFont=*-italic,BoldItalicFont=*-bolditalic]
  \setmonofont{DejaVuSansMono}[Extension=.ttf,UprightFont=*,BoldFont=*-Bold,ItalicFont=*-Oblique,BoldItalicFont=*-BoldOblique,Scale=0.86]
\fi
\usepackage[ngerman,shorthands=off]{babel}
\usepackage{amssymb}
\usepackage{xcolor,graphicx,tabularx,array,enumitem,fancyhdr,tikz,multicol,calc,etoolbox,pifont,xspace,listings}
\usepackage[most]{tcolorbox}
\usetikzlibrary{arrows.meta,positioning,calc,shapes.geometric,decorations.pathreplacing}
\usepackage[hidelinks]{hyperref}
\setlength{\parindent}{0pt}
\setlength{\parskip}{3pt}
\setlist{nosep,leftmargin=*}

% ---------- Lösungsschalter ----------
\newif\ifloesung
\ifdefined\mitloesung\loesungtrue\fi
\newcommand{\lsgfarbe}{\color{red!75!black}}

% ---------- Kopf- und Fußzeile ----------
\newcommand{\lehrkraft}{Hau}
\newcommand{\themenbereich}{Künstliche Intelligenz}
\newcommand{\abnummer}{}
\pagestyle{fancy}\fancyhf{}
\renewcommand{\headrulewidth}{0.4pt}
\fancyhead[L]{\small Klasse 13\\Datum:}
\fancyhead[C]{\small Themenbereich:\\\themenbereich}
\fancyhead[R]{\small \lehrkraft\ifloesung\\{\lsgfarbe\bfseries Lösung}\fi}
\fancyfoot[C]{\scriptsize edu-mrh.de\,·\,Informatik 13\,·\,Blatt \abnummer\ifloesung\ (Lösung)\fi\,·\,Seite \thepage}
\newcommand{\abtitel}[2]{\renewcommand{\abnummer}{#1}\begin{center}\Large\bfseries #1\quad\underline{#2}\end{center}\vspace{-2mm}}

% ---------- Boxen ----------
\newenvironment{aufgabe}[2][]{\begin{tcolorbox}[enhanced,breakable,colback=white,colframe=black,boxrule=0.7pt,arc=3mm,left=2mm,right=2mm,top=1mm,bottom=1.2mm,before skip=2mm,after skip=2mm]\textbf{Aufgabe #2}\ifstrempty{#1}{}{ \textit{#1}}\par\vspace{0.3mm}}{\end{tcolorbox}}
\newtcolorbox{merke}[1][Merke]{enhanced,breakable,colback=green!5,colframe=green!40!black,boxrule=0.8pt,arc=2mm,left=2mm,right=2mm,top=1mm,bottom=1.2mm,before skip=2mm,after skip=2mm,title={#1},fonttitle=\bfseries,coltitle=white,toptitle=0.3mm,bottomtitle=0.3mm}
\newtcolorbox{hinweis}{enhanced,breakable,colback=yellow!12,colframe=orange!70!black,boxrule=0.6pt,arc=2mm,left=2mm,right=2mm,top=1mm,bottom=1mm,before skip=2mm,after skip=2mm}
\newcommand{\web}[1]{{\small\color{gray}\textit{Website: #1}}}
\newcommand{\code}[1]{\texttt{#1}}

% ---------- Antwortplatz (wie informatik/10/java, Lücken x 1,5) ----------
\newlength{\zeilenabstand}\setlength{\zeilenabstand}{8mm}
\newlength{\karo}\setlength{\karo}{5mm}
\tikzset{karo/.style={draw=gray!50,line width=0.3pt},schreiblinie/.style={draw=black!65,line width=0.35pt}}
\newlength{\lsgbreite}
\newcommand{\lsg}[2][6.75cm]{\leavevmode\ifloesung\settowidth{\lsgbreite}{\,#2\,}%
  \ifdim\lsgbreite>#1\relax\rlap{\raisebox{0.8pt}{\resizebox{#1}{!}{\lsgfarbe\,#2\,}}}%
  \else\rlap{\raisebox{0.8pt}{\lsgfarbe\,#2}}\fi\fi\rule[-0.6ex]{#1}{0.4pt}\xspace}
\newcommand{\luecke}[1][6.75cm]{\lsg[#1]{}}
\newcommand{\lsgrest}[1]{\leavevmode\unskip\hspace{1.5mm}\ifloesung\rlap{\raisebox{0.8pt}{\lsgfarbe\,#1}}\fi\leaders\hbox{\rule[-0.6ex]{0.5pt}{0.4pt}}\hfill\null\par}
\newcommand{\restzeile}{\lsgrest{}}
\newcommand{\lsgtext}[1]{\ifloesung\leavevmode{\lsgfarbe#1}\fi}
\newcommand{\lsgoder}[2]{\ifloesung\leavevmode#1\else#2\fi}
\newenvironment{lueckentext}{\par\linespread{1.5}\selectfont\raggedright}{\par}
\newcommand{\schreibzeile}{\rule[-2.8mm]{0pt}{8mm}}
% Tabellenzelle für n Zeilen Handschrift (Lösungstext darin kleiner setzen: \lsgklein{…})
\newcommand{\schreibzeilen}[1]{\rule[\dimexpr-#1\zeilenabstand+5.2mm\relax]{0pt}{\dimexpr#1\zeilenabstand\relax}}
\newcommand{\lsgklein}[1]{\lsgtext{\footnotesize\raggedright #1}}
\newcommand{\nr}[1]{\ding{\numexpr171+#1\relax}}
\newcommand{\lsglinien}[2]{\par\vspace{0.5mm}\noindent\begin{tikzpicture}
  \path (0,0) (\linewidth-0.4pt,-#1\zeilenabstand-1.5mm);
  \foreach \i in {1,...,#1}{\draw[schreiblinie] (0.2pt,-\i\zeilenabstand) -- (\linewidth-0.2pt,-\i\zeilenabstand);}
  \ifloesung\node[overlay,anchor=north west,inner sep=0pt] at (1mm,-\zeilenabstand+1pt+5mm) {\parbox[t]{\linewidth-2mm}{\lsgfarbe\raggedright\setlength{\baselineskip}{\zeilenabstand}\rule{0pt}{5mm}#2}};\fi
  \end{tikzpicture}\par}
\newcommand{\linien}[1]{\lsglinien{#1}{}}
\newcommand{\lsgfeld}[2]{\noindent\begin{tikzpicture}
  \draw[line width=0.4pt] (0.2pt,0) rectangle (\linewidth-0.2pt,-#1\zeilenabstand-3mm);
  \foreach \i in {1,...,#1}{\draw[schreiblinie] (2mm,-\i\zeilenabstand) -- (\linewidth-2mm,-\i\zeilenabstand);}
  \ifloesung\node[overlay,anchor=north west,inner sep=0pt] at (2.5mm,-\zeilenabstand+1pt+5mm) {\parbox[t]{\linewidth-5mm}{\lsgfarbe\ttfamily\small\setlength{\baselineskip}{\zeilenabstand}\rule{0pt}{5mm}#2}};\fi
  \end{tikzpicture}}
\newcommand{\linienfeld}[1]{\lsgfeld{#1}{}}
\newcommand{\karoraster}[3]{\begin{tikzpicture}\draw[karo] (0,0) grid[step=\karo] (#1\karo,#2\karo);
  \ifloesung\begin{scope}[overlay,shift={(0,#2\karo)},red!75!black,text=red!75!black,line width=0.8pt]#3\end{scope}\fi\end{tikzpicture}}
\newcommand{\lsgkariert}[2]{\par\vspace{1.5mm}\noindent\pgfmathtruncatemacro{\karoX}{floor(\linewidth/\karo)}%
  \makebox[\linewidth]{\karoraster{\karoX}{#1}{#2}}\par\vspace{1mm}}
\newcommand{\kariert}[1]{\lsgkariert{#1}{}}
\newcommand{\lsgkarobox}[3]{\pgfmathtruncatemacro{\karoX}{round((#1)/\karo)}\pgfmathtruncatemacro{\karoY}{round((#2)/\karo)}\karoraster{\karoX}{\karoY}{#3}}
\newcommand{\karobox}[2]{\lsgkarobox{#1}{#2}{}}
% Ankreuzfeld: \kreuz{1} = angekreuzt (nur in der Lösung)
\newcommand{\kreuz}[1]{\leavevmode\ifloesung\ifnum#1=1 {\lsgfarbe$\boxtimes$}\else$\square$\fi\else$\square$\fi}

% ---------- Listings: Java und Python ----------
\lstset{basicstyle=\ttfamily\small,keywordstyle=\bfseries,commentstyle=\color{gray!80!black},stringstyle=\color{black},showstringspaces=false,columns=flexible,keepspaces=true,tabsize=3,frame=single,framerule=0.4pt,rulecolor=\color{black},xleftmargin=2mm,framexleftmargin=1mm,aboveskip=1.5mm,belowskip=1.5mm,breaklines=true,escapeinside={(*}{*)},
 literate={ä}{{\"a}}1 {ö}{{\"o}}1 {ü}{{\"u}}1 {Ä}{{\"A}}1 {Ö}{{\"O}}1 {Ü}{{\"U}}1 {ß}{{\ss}}1 {–}{{--}}1 {…}{{\dots}}1 {≤}{{$\le$}}1 {≥}{{$\ge$}}1 {≠}{{$\ne$}}1 {·}{{$\cdot$}}1 {①}{{\nr{1}}}1 {②}{{\nr{2}}}1 {③}{{\nr{3}}}1 {④}{{\nr{4}}}1}
\lstnewenvironment{java}[1][]{\lstset{language=Java,#1}}{}
\lstnewenvironment{javaluecke}[1][]{\lstset{language=Java,basicstyle=\linespread{1.5}\small\ttfamily,#1}}{}
\lstnewenvironment{python}[1][]{\lstset{language=Python,#1}}{}
\lstnewenvironment{pythonluecke}[1][]{\lstset{language=Python,basicstyle=\linespread{1.5}\small\ttfamily,#1}}{}
\lstnewenvironment{prolog}[1][]{\lstset{morecomment=[l]{\%},#1}}{}
\lstnewenvironment{prologluecke}[1][]{\lstset{morecomment=[l]{\%},basicstyle=\linespread{1.5}\small\ttfamily,#1}}{}

\newcommand{\niveauEA}{\tikz[baseline=(t.base)]\node[fill=orange!20,draw=orange!70!black,rounded corners=1.5mm,inner sep=1pt,font=\scriptsize\bfseries,text=orange!50!black](t){\,eA\,};}

% ---------- TikZ-Stile ----------
\tikzset{
  neuron/.style={circle,draw=black,line width=0.7pt,minimum size=8mm,inner sep=0pt,font=\small},
  eingabe/.style={neuron,fill=gray!12},
  kante/.style={-{Stealth[length=2mm]},line width=0.6pt},
  gew/.style={font=\scriptsize,fill=white,inner sep=0.6pt},
  pt/.style={circle,inner sep=0pt,minimum size=2.2mm,draw=white,line width=0.3pt},
}
% ---------- Farben wie auf der Website ----------
\definecolor{adelie}{HTML}{EA580C}\definecolor{zuegel}{HTML}{9333EA}\definecolor{esels}{HTML}{0E7490}
\definecolor{klnull}{HTML}{2563EB}\definecolor{kleins}{HTML}{DC2626}
\definecolor{cA}{HTML}{2563EB}\definecolor{cB}{HTML}{DC2626}\definecolor{cC}{HTML}{16A34A}\definecolor{cD}{HTML}{D97706}\definecolor{cE}{HTML}{9333EA}\definecolor{cF}{HTML}{0891B2}
\setlength{\unitlength}{1mm}

\begin{document}
\abtitel{3}{k-Means: Ordnung ohne Lehrer}

\begin{aufgabe}[k-Means von Hand, $k=3$]{1}
\begin{minipage}[t]{0.6\linewidth}\vspace{0pt}
\begin{tikzpicture}[x=5mm,y=5mm]
  \draw[karo] (0,0) grid (19,19);
  \draw[-{Stealth[length=2mm]},line width=0.6pt] (0,0) -- (19.6,0) node[below left,font=\footnotesize]{$x$};
  \draw[-{Stealth[length=2mm]},line width=0.6pt] (0,0) -- (0,19.6) node[below left,font=\footnotesize]{$y$};
  \foreach \v in {1,...,9}{\node[font=\tiny,below] at (2*\v,0) {\v}; \node[font=\tiny,left] at (0,2*\v) {\v};}
  \foreach \n/\x/\y in {A/1/2,B/2/1,C/2/3,D/3/2,E/7/8,F/8/7,G/8/9,H/9/8,I/6/2,J/8/1,K/9/3,L/5/6}{
    \fill (2*\x,2*\y) circle (2.4pt); \node[font=\scriptsize,above right,inner sep=1pt] at (2*\x,2*\y) {\n};}
  \ifloesung
    \begin{scope}[red!75!black,line width=0.9pt]
      \foreach \x/\y/\t in {2/3/1,9/8/2,8/1/3}{\draw (2*\x-0.5,2*\y-0.5) -- (2*\x+0.5,2*\y+0.5) (2*\x-0.5,2*\y+0.5) -- (2*\x+0.5,2*\y-0.5);}
      \draw[dashed,-{Stealth}] (4,6) -- (5.2,5.6); \draw[dashed,-{Stealth}] (5.2,5.6) -- (4,4);
      \draw[dashed,-{Stealth}] (18,16) -- (16,16); \draw[dashed,-{Stealth}] (16,16) -- (14.8,15.2);
      \draw[dashed,-{Stealth}] (16,2) -- (15.33,4);
      \foreach \x/\y in {4/4,14.8/15.2,15.33/4}{\draw[line width=1.3pt] (\x-0.55,\y-0.55) -- (\x+0.55,\y+0.55) (\x-0.55,\y+0.55) -- (\x+0.55,\y-0.55);}
    \end{scope}
  \fi
\end{tikzpicture}
\end{minipage}\hfill
\begin{minipage}[t]{0.38\linewidth}\vspace{0pt}
Startzentren: Z1 = C(2\,\textbar\,3), Z2 = H(9\,\textbar\,8), Z3 = J(8\,\textbar\,1).\par\vspace{1mm}
\nr{1}~Jeder Punkt geht zum nächsten Zentrum.\\
\nr{2}~Jedes Zentrum wandert in die Mitte seiner Punkte.\\
\nr{3}~Wiederholen, bis sich nichts mehr ändert.\par\vspace{1.5mm}
Zeichne die Zentren als Kreuz ein und ihre Wege als Pfeile.\par\vspace{1.5mm}
Nur bei L ist es knapp. Berechne im 2.~Durchgang die Abstände von L zu Z1 und Z2:
\begin{lueckentext}
$d(L,Z1) = $ \lsg[4.2cm]{$\sqrt{2{,}4^2+3{,}2^2}=4$}\\
$d(L,Z2) = $ \lsg[4.2cm]{$\sqrt{3^2+2^2}\approx 3{,}61$}\\
L gehört dann zu \lsg[2.25cm]{Z2}.
\end{lueckentext}
\end{minipage}\par\vspace{2mm}
\renewcommand{\arraystretch}{1.1}
\begin{tabularx}{\linewidth}{|c|X|X|X|p{4.7cm}|}\hline
\textbf{Durchgang} & \textbf{Punkte bei Z1} & \textbf{Punkte bei Z2} & \textbf{Punkte bei Z3} & \textbf{neue Zentren (\nr{2})} \\\hline
1\schreibzeilen{2} & \lsgklein{A, B, C, D, L} & \lsgklein{E, F, G, H} & \lsgklein{I, J, K} & \lsgklein{Z1(2,6\,|\,2,8), Z2(8\,|\,8),\\ Z3(7,67\,|\,2)} \\\hline
2\schreibzeilen{2} & \lsgklein{A, B, C, D} & \lsgklein{E, F, G, H, L} & \lsgklein{I, J, K} & \lsgklein{Z1(2\,|\,2), Z2(7,4\,|\,7,6),\\ Z3(7,67\,|\,2)} \\\hline
3\schreibzeilen{1} & \lsgklein{A, B, C, D} & \lsgklein{E, F, G, H, L} & \lsgklein{I, J, K} & \lsgklein{keine Änderung: fertig} \\\hline
\end{tabularx}
\end{aufgabe}

\begin{aufgabe}[Der Algorithmus (Website: Festhalten)]{2}
\begin{lueckentext}
k-Means ist ein Verfahren des \lsg[5.25cm]{unüberwachten} Lernens: Die Punkte haben keine \lsg[5.25cm]{Etiketten (Klassen)}. Vorgegeben wird nur die Zahl \lsg[2.25cm]{$k$} der Cluster. \nr{1}~Jeder Punkt wird dem Zentrum mit dem kleinsten \lsg[4.5cm]{Abstand} zugeordnet. \nr{2}~Jedes Zentrum wird auf den \lsg[5.25cm]{Mittelwert} seiner Punkte gesetzt. \nr{3}~Beides wird wiederholt, bis sich \lsgrest{beim Zuordnen nichts mehr ändert.}
Die Streuung (Summe der quadrierten Abstände zum eigenen Zentrum) wird dabei nie \lsg[4.5cm]{größer}.
\end{lueckentext}
\end{aufgabe}

\newpage
\begin{aufgabe}[Selbst bauen (Online-IDE, TODO 3b und 3d)]{3}
Ergänze die Lücken in den Kernzeilen von \code{zuordnen()} und \code{lauf()}.
\begin{javaluecke}
for (int i = 0; i < punkte.length; i++) {                    // (*\nr{1}*) zuordnen()
   int bester = 0;  double kleinster = abstand(punkte[i].getWerte(), zentren[0]);
   for (int c = 1; c < k; c++) {
      double d = abstand(punkte[i].getWerte(), (*\lsg[2.8cm]{zentren[c]}*));
      if ((*\lsg[3.6cm]{d < kleinster}*)) { kleinster = d;  bester = c; }
   }
   if (punkte[i].getCluster() != bester) { (*\lsg[5.2cm]{punkte[i].setCluster(bester);}*) geaendert = true; }
}
while ((*\lsg[2.8cm]{zuordnen()}*)) { (*\lsg[4.4cm]{zentrenVerschieben();}*) durchgaenge++; }     // (*\nr{3}*) lauf()
\end{javaluecke}
\end{aufgabe}

\begin{aufgabe}[Startpunkte (Online-IDE)]{4}
\renewcommand{\arraystretch}{1.1}
\begin{tabularx}{\linewidth}{|p{2.8cm}|p{1.9cm}|p{1.9cm}|X|}\hline
\textbf{Start} & \textbf{Durchg.} & \textbf{Streuung} & \textbf{Was fällt auf?} \\\hline
\{0, 200, 300\}\schreibzeilen{2} & \lsgklein{3} & \lsgklein{13,55} & \lsgklein{Esels fast perfekt in einem Cluster; Adélie und Zügel nur teilweise getrennt (11 Adélie bei den Zügeln).} \\\hline
\{150, 160, 170\}\schreibzeilen{2} & \lsgklein{5} & \lsgklein{18,41} & \lsgklein{Esels auf zwei Cluster geteilt, Adélie und Zügel zusammen: schlechtere Lösung (lokales Minimum).} \\\hline
\end{tabularx}\par\vspace{1mm}
Wie kann man vermeiden, dass k-Means in einer schlechten Lösung hängen bleibt?
\lsglinien{1}{Mehrmals mit verschiedenen Starts rechnen, das Ergebnis mit der kleinsten Streuung nehmen.}
\end{aufgabe}

\begin{aufgabe}[Wie viele Cluster? Der Ellbogen (Online-IDE, TODO 5)]{5}
\begin{minipage}[t]{0.5\linewidth}\vspace{0pt}
\begin{tikzpicture}[x=5mm,y=5mm]
  \draw[karo] (0,0) grid (14,10);
  \draw[-{Stealth[length=2mm]},line width=0.6pt] (0,0) -- (14.6,0) node[below left,font=\footnotesize]{$k$};
  \draw[-{Stealth[length=2mm]},line width=0.6pt] (0,0) -- (0,10.6) node[right,font=\footnotesize]{Streuung};
  \foreach \v in {1,...,6}{\node[font=\tiny,below] at (2*\v,0) {\v};}
  \foreach \v in {10,20,...,50}{\node[font=\tiny,left] at (0,\v/5.5) {\v};}
  \ifloesung
    \draw[red!75!black,line width=0.9pt] (2,9.39) -- (4,3.81) -- (6,2.46) -- (8,2.01) -- (10,1.59) -- (12,1.36);
    \foreach \x/\y in {2/9.39,4/3.81,6/2.46,8/2.01,10/1.59,12/1.36}{\fill[red!75!black] (\x,\y) circle (2.2pt);}
  \fi
\end{tikzpicture}
\end{minipage}\hfill
\begin{minipage}[t]{0.47\linewidth}\vspace{0pt}
Trage die Streuung für $k = 1$ bis 6 ein und verbinde die Punkte.
\begin{lueckentext}
Der Knick liegt bei $k = $ \lsg[2.25cm]{3}.
Begründung: \lsgrest{Danach bringt jedes}
\lsgrest{weitere Cluster nur noch wenig.}
\end{lueckentext}
Warum ist das $k$ mit der kleinsten Streuung nicht das beste?
\lsglinien{2}{Sie sinkt mit jedem Cluster; bei k = 342 wäre sie 0 (jeder Pinguin allein): nutzlos.}
\end{minipage}
\end{aufgabe}

\begin{aufgabe}[Grenzen (Website)]{6}
Bei den „zwei Monden“ und beim „Ring um Kern“ findet k-Means nicht die Gruppen, die ein Mensch sieht. Erkläre mit der Lage der Grenze zwischen zwei Zentren.
\lsglinien{2}{Die Grenze zwischen zwei Zentren ist eine Gerade (Mittelsenkrechte). Gebogene oder ineinanderliegende Gruppen trennt sie nicht; k-Means passt zu runden Gruppen.}
\end{aufgabe}
\end{document}
)
%  Lösungseinträge belegen nie mehr Platz als die Lücke im Blatt.
%
%  Antwortplatz:
%   \luecke[l] / \lsg[l]{Text}      Lücke (Standard 6,75 cm, Fachbegriff 5,25 cm, Zahl 2,25 cm)
%   \restzeile / \lsgrest{Text}     Lücke bis Zeilenende      \lueckentext  Absatz mit Lücken
%   \linien{n} / \lsglinien{n}{T}   Schreiblinien 8 mm        \linienfeld{n} / \lsgfeld{n}{Code}
%   \schreibzeile                   Tabellenzelle 8 mm        \lsgtext{…}  nur in der Lösung
%   \kariert{n} / \lsgkariert{n}{TikZ}, \karobox{b}{h} / \lsgkarobox{b}{h}{TikZ}
%   \kreuz{0|1}                     Ankreuzfeld (1 = in der Lösung angekreuzt)
%  KI:
%   java / javaluecke               Listings (Lücken mit (*\lsg[3cm]{…}*))
%   \nr{1}                          Schrittnummer ① wie auf der Website und in der IDE
%   \niveauEA                       Kennzeichen „eA“ für Aufgaben des erhöhten Niveaus
%   Farben: adelie, zuegel, esels (Pinguinarten), klnull, kleins (Klassen 0/1), cA … cF (Cluster)
% =====================================================================
\usepackage[a4paper,left=16mm,right=16mm,top=23mm,bottom=13mm,headheight=12mm,headsep=3mm,footskip=6mm]{geometry}
\usepackage{iftex}
\ifPDFTeX
  \usepackage[T1]{fontenc}
  \usepackage[utf8]{inputenc}
  \usepackage{helvet}
  \usepackage[scaled=0.86]{beramono}
  \renewcommand{\familydefault}{\sfdefault}
\else
  \usepackage{fontspec}
  \setmainfont{texgyreheros}[Extension=.otf,UprightFont=*-regular,BoldFont=*-bold,ItalicFont=*-italic,BoldItalicFont=*-bolditalic]
  \setmonofont{DejaVuSansMono}[Extension=.ttf,UprightFont=*,BoldFont=*-Bold,ItalicFont=*-Oblique,BoldItalicFont=*-BoldOblique,Scale=0.86]
\fi
\usepackage[ngerman,shorthands=off]{babel}
\usepackage{amssymb}
\usepackage{xcolor,graphicx,tabularx,array,enumitem,fancyhdr,tikz,multicol,calc,etoolbox,pifont,xspace,listings}
\usepackage[most]{tcolorbox}
\usetikzlibrary{arrows.meta,positioning,calc,shapes.geometric,decorations.pathreplacing}
\usepackage[hidelinks]{hyperref}
\setlength{\parindent}{0pt}
\setlength{\parskip}{3pt}
\setlist{nosep,leftmargin=*}

% ---------- Lösungsschalter ----------
\newif\ifloesung
\ifdefined\mitloesung\loesungtrue\fi
\newcommand{\lsgfarbe}{\color{red!75!black}}

% ---------- Kopf- und Fußzeile ----------
\newcommand{\lehrkraft}{Hau}
\newcommand{\themenbereich}{Künstliche Intelligenz}
\newcommand{\abnummer}{}
\pagestyle{fancy}\fancyhf{}
\renewcommand{\headrulewidth}{0.4pt}
\fancyhead[L]{\small Klasse 13\\Datum:}
\fancyhead[C]{\small Themenbereich:\\\themenbereich}
\fancyhead[R]{\small \lehrkraft\ifloesung\\{\lsgfarbe\bfseries Lösung}\fi}
\fancyfoot[C]{\scriptsize edu-mrh.de\,·\,Informatik 13\,·\,Blatt \abnummer\ifloesung\ (Lösung)\fi\,·\,Seite \thepage}
\newcommand{\abtitel}[2]{\renewcommand{\abnummer}{#1}\begin{center}\Large\bfseries #1\quad\underline{#2}\end{center}\vspace{-2mm}}

% ---------- Boxen ----------
\newenvironment{aufgabe}[2][]{\begin{tcolorbox}[enhanced,breakable,colback=white,colframe=black,boxrule=0.7pt,arc=3mm,left=2mm,right=2mm,top=1mm,bottom=1.2mm,before skip=2mm,after skip=2mm]\textbf{Aufgabe #2}\ifstrempty{#1}{}{ \textit{#1}}\par\vspace{0.3mm}}{\end{tcolorbox}}
\newtcolorbox{merke}[1][Merke]{enhanced,breakable,colback=green!5,colframe=green!40!black,boxrule=0.8pt,arc=2mm,left=2mm,right=2mm,top=1mm,bottom=1.2mm,before skip=2mm,after skip=2mm,title={#1},fonttitle=\bfseries,coltitle=white,toptitle=0.3mm,bottomtitle=0.3mm}
\newtcolorbox{hinweis}{enhanced,breakable,colback=yellow!12,colframe=orange!70!black,boxrule=0.6pt,arc=2mm,left=2mm,right=2mm,top=1mm,bottom=1mm,before skip=2mm,after skip=2mm}
\newcommand{\web}[1]{{\small\color{gray}\textit{Website: #1}}}
\newcommand{\code}[1]{\texttt{#1}}

% ---------- Antwortplatz (wie informatik/10/java, Lücken x 1,5) ----------
\newlength{\zeilenabstand}\setlength{\zeilenabstand}{8mm}
\newlength{\karo}\setlength{\karo}{5mm}
\tikzset{karo/.style={draw=gray!50,line width=0.3pt},schreiblinie/.style={draw=black!65,line width=0.35pt}}
\newlength{\lsgbreite}
\newcommand{\lsg}[2][6.75cm]{\leavevmode\ifloesung\settowidth{\lsgbreite}{\,#2\,}%
  \ifdim\lsgbreite>#1\relax\rlap{\raisebox{0.8pt}{\resizebox{#1}{!}{\lsgfarbe\,#2\,}}}%
  \else\rlap{\raisebox{0.8pt}{\lsgfarbe\,#2}}\fi\fi\rule[-0.6ex]{#1}{0.4pt}\xspace}
\newcommand{\luecke}[1][6.75cm]{\lsg[#1]{}}
\newcommand{\lsgrest}[1]{\leavevmode\unskip\hspace{1.5mm}\ifloesung\rlap{\raisebox{0.8pt}{\lsgfarbe\,#1}}\fi\leaders\hbox{\rule[-0.6ex]{0.5pt}{0.4pt}}\hfill\null\par}
\newcommand{\restzeile}{\lsgrest{}}
\newcommand{\lsgtext}[1]{\ifloesung\leavevmode{\lsgfarbe#1}\fi}
\newcommand{\lsgoder}[2]{\ifloesung\leavevmode#1\else#2\fi}
\newenvironment{lueckentext}{\par\linespread{1.5}\selectfont\raggedright}{\par}
\newcommand{\schreibzeile}{\rule[-2.8mm]{0pt}{8mm}}
% Tabellenzelle für n Zeilen Handschrift (Lösungstext darin kleiner setzen: \lsgklein{…})
\newcommand{\schreibzeilen}[1]{\rule[\dimexpr-#1\zeilenabstand+5.2mm\relax]{0pt}{\dimexpr#1\zeilenabstand\relax}}
\newcommand{\lsgklein}[1]{\lsgtext{\footnotesize\raggedright #1}}
\newcommand{\nr}[1]{\ding{\numexpr171+#1\relax}}
\newcommand{\lsglinien}[2]{\par\vspace{0.5mm}\noindent\begin{tikzpicture}
  \path (0,0) (\linewidth-0.4pt,-#1\zeilenabstand-1.5mm);
  \foreach \i in {1,...,#1}{\draw[schreiblinie] (0.2pt,-\i\zeilenabstand) -- (\linewidth-0.2pt,-\i\zeilenabstand);}
  \ifloesung\node[overlay,anchor=north west,inner sep=0pt] at (1mm,-\zeilenabstand+1pt+5mm) {\parbox[t]{\linewidth-2mm}{\lsgfarbe\raggedright\setlength{\baselineskip}{\zeilenabstand}\rule{0pt}{5mm}#2}};\fi
  \end{tikzpicture}\par}
\newcommand{\linien}[1]{\lsglinien{#1}{}}
\newcommand{\lsgfeld}[2]{\noindent\begin{tikzpicture}
  \draw[line width=0.4pt] (0.2pt,0) rectangle (\linewidth-0.2pt,-#1\zeilenabstand-3mm);
  \foreach \i in {1,...,#1}{\draw[schreiblinie] (2mm,-\i\zeilenabstand) -- (\linewidth-2mm,-\i\zeilenabstand);}
  \ifloesung\node[overlay,anchor=north west,inner sep=0pt] at (2.5mm,-\zeilenabstand+1pt+5mm) {\parbox[t]{\linewidth-5mm}{\lsgfarbe\ttfamily\small\setlength{\baselineskip}{\zeilenabstand}\rule{0pt}{5mm}#2}};\fi
  \end{tikzpicture}}
\newcommand{\linienfeld}[1]{\lsgfeld{#1}{}}
\newcommand{\karoraster}[3]{\begin{tikzpicture}\draw[karo] (0,0) grid[step=\karo] (#1\karo,#2\karo);
  \ifloesung\begin{scope}[overlay,shift={(0,#2\karo)},red!75!black,text=red!75!black,line width=0.8pt]#3\end{scope}\fi\end{tikzpicture}}
\newcommand{\lsgkariert}[2]{\par\vspace{1.5mm}\noindent\pgfmathtruncatemacro{\karoX}{floor(\linewidth/\karo)}%
  \makebox[\linewidth]{\karoraster{\karoX}{#1}{#2}}\par\vspace{1mm}}
\newcommand{\kariert}[1]{\lsgkariert{#1}{}}
\newcommand{\lsgkarobox}[3]{\pgfmathtruncatemacro{\karoX}{round((#1)/\karo)}\pgfmathtruncatemacro{\karoY}{round((#2)/\karo)}\karoraster{\karoX}{\karoY}{#3}}
\newcommand{\karobox}[2]{\lsgkarobox{#1}{#2}{}}
% Ankreuzfeld: \kreuz{1} = angekreuzt (nur in der Lösung)
\newcommand{\kreuz}[1]{\leavevmode\ifloesung\ifnum#1=1 {\lsgfarbe$\boxtimes$}\else$\square$\fi\else$\square$\fi}

% ---------- Listings: Java und Python ----------
\lstset{basicstyle=\ttfamily\small,keywordstyle=\bfseries,commentstyle=\color{gray!80!black},stringstyle=\color{black},showstringspaces=false,columns=flexible,keepspaces=true,tabsize=3,frame=single,framerule=0.4pt,rulecolor=\color{black},xleftmargin=2mm,framexleftmargin=1mm,aboveskip=1.5mm,belowskip=1.5mm,breaklines=true,escapeinside={(*}{*)},
 literate={ä}{{\"a}}1 {ö}{{\"o}}1 {ü}{{\"u}}1 {Ä}{{\"A}}1 {Ö}{{\"O}}1 {Ü}{{\"U}}1 {ß}{{\ss}}1 {–}{{--}}1 {…}{{\dots}}1 {≤}{{$\le$}}1 {≥}{{$\ge$}}1 {≠}{{$\ne$}}1 {·}{{$\cdot$}}1 {①}{{\nr{1}}}1 {②}{{\nr{2}}}1 {③}{{\nr{3}}}1 {④}{{\nr{4}}}1}
\lstnewenvironment{java}[1][]{\lstset{language=Java,#1}}{}
\lstnewenvironment{javaluecke}[1][]{\lstset{language=Java,basicstyle=\linespread{1.5}\small\ttfamily,#1}}{}
\lstnewenvironment{python}[1][]{\lstset{language=Python,#1}}{}
\lstnewenvironment{pythonluecke}[1][]{\lstset{language=Python,basicstyle=\linespread{1.5}\small\ttfamily,#1}}{}
\lstnewenvironment{prolog}[1][]{\lstset{morecomment=[l]{\%},#1}}{}
\lstnewenvironment{prologluecke}[1][]{\lstset{morecomment=[l]{\%},basicstyle=\linespread{1.5}\small\ttfamily,#1}}{}

\newcommand{\niveauEA}{\tikz[baseline=(t.base)]\node[fill=orange!20,draw=orange!70!black,rounded corners=1.5mm,inner sep=1pt,font=\scriptsize\bfseries,text=orange!50!black](t){\,eA\,};}

% ---------- TikZ-Stile ----------
\tikzset{
  neuron/.style={circle,draw=black,line width=0.7pt,minimum size=8mm,inner sep=0pt,font=\small},
  eingabe/.style={neuron,fill=gray!12},
  kante/.style={-{Stealth[length=2mm]},line width=0.6pt},
  gew/.style={font=\scriptsize,fill=white,inner sep=0.6pt},
  pt/.style={circle,inner sep=0pt,minimum size=2.2mm,draw=white,line width=0.3pt},
}
% ---------- Farben wie auf der Website ----------
\definecolor{adelie}{HTML}{EA580C}\definecolor{zuegel}{HTML}{9333EA}\definecolor{esels}{HTML}{0E7490}
\definecolor{klnull}{HTML}{2563EB}\definecolor{kleins}{HTML}{DC2626}
\definecolor{cA}{HTML}{2563EB}\definecolor{cB}{HTML}{DC2626}\definecolor{cC}{HTML}{16A34A}\definecolor{cD}{HTML}{D97706}\definecolor{cE}{HTML}{9333EA}\definecolor{cF}{HTML}{0891B2}
\setlength{\unitlength}{1mm}

\begin{document}
\abtitel{W3}{Die Inferenzmaschine \niveauEA}

\begin{aufgabe}[Selbst Inferenzmaschine sein]{1}
\begin{prolog}
1  elternteil(oskar, pia).     4  vorfahre(X, Y) :- elternteil(X, Y).
2  elternteil(pia, tim).       5  vorfahre(X, Y) :- elternteil(X, Z), vorfahre(Z, Y).
3  elternteil(tim, kim).
\end{prolog}
Anfrage \code{?- vorfahre(oskar, kim).} Notiere die ersten Schritte.\par\vspace{1mm}
\renewcommand{\arraystretch}{1.1}
\begin{tabularx}{\linewidth}{|>{\raggedright\arraybackslash}p{4.2cm}|c|X|c|}\hline
\textbf{Ziel} & \textbf{Klausel} & \textbf{Belegung / neue Ziele} & \textbf{Ergebnis} \\\hline
\code{vorfahre(oskar, kim)}\schreibzeile & \lsgklein{4} & \lsgklein{X = oskar, Y = kim; neues Ziel elternteil(oskar, kim)} & \lsgklein{weiter} \\\hline
\lsgklein{elternteil(oskar, kim)}\schreibzeile & \lsgklein{1–3} & \lsgklein{kein Fakt passt} & \lsgklein{\ding{55} zurück} \\\hline
\lsgklein{vorfahre(oskar, kim)}\schreibzeile & \lsgklein{5} & \lsgklein{neue Ziele elternteil(oskar, Z), vorfahre(Z, kim)} & \lsgklein{weiter} \\\hline
\lsgklein{elternteil(oskar, Z)}\schreibzeile & \lsgklein{1} & \lsgklein{Z = pia} & \lsgklein{\ding{51}} \\\hline
\lsgklein{vorfahre(pia, kim)}\schreibzeile & \lsgklein{4, dann 5} & \lsgklein{wie oben: über Z = tim zu vorfahre(tim, kim), Klausel 4 mit Fakt 3} & \lsgklein{ja} \\\hline
\end{tabularx}
\end{aufgabe}

\begin{aufgabe}[Der Ableitungsbaum]{2}
Zeichne den Ableitungsbaum für \code{?- vorfahre(oskar, kim).}: Knoten = Ziele, Kanten = verwendete Klausel. Markiere Sackgassen mit \ding{55} und die Antwort mit $\bigstar$.
\lsgkariert{11}{
  \begin{scope}[font=\scriptsize,yscale=-1,inner sep=1.5pt]
  \node[draw,fill=white] (a) at (8.5,0.5) {vorfahre(oskar, kim)};
  \node[draw,fill=white] (b) at (3.5,1.3) {elternteil(oskar, kim) \ding{55}};
  \node[draw,fill=white] (c) at (12,1.3) {elternteil(oskar, Z), vorfahre(Z, kim)};
  \node[draw,fill=white] (d) at (12,2.1) {vorfahre(pia, kim)};
  \node[draw,fill=white] (e) at (7.5,2.9) {elternteil(pia, kim) \ding{55}};
  \node[draw,fill=white] (f) at (13.2,2.9) {elternteil(pia, Z2), vorfahre(Z2, kim)};
  \node[draw,fill=white] (g) at (13.2,3.7) {vorfahre(tim, kim)};
  \node[draw,fill=white] (hh) at (10,4.6) {elternteil(tim, kim) $\bigstar$};
  \draw (a) -- node[above] {4} (b); \draw (a) -- node[above] {5} (c); \draw (c) -- node[right] {1: Z = pia} (d);
  \draw (d) -- node[above] {4} (e); \draw (d) -- node[right] {5} (f); \draw (f) -- node[right] {2: Z2 = tim} (g); \draw (g) -- node[above] {4} (hh);
  \end{scope}}
\end{aufgabe}

\begin{aufgabe}[Merkkasten]{3}
\begin{lueckentext}
Die Inferenzmaschine bearbeitet immer das \lsg[2.6cm]{erste} Ziel und sucht die Klauseln von \lsg[4.2cm]{oben nach unten} ab. Ziel und Klauselkopf werden durch \lsg[4.2cm]{Unifikation} gleich gemacht, dabei werden Variablen belegt. Bei einer Regel wird das Ziel durch den \lsg[2.6cm]{Rumpf} ersetzt. Scheitert ein Ziel, geht sie zur letzten Wahlmöglichkeit zurück: \lsg[4.2cm]{Backtracking}. Den Ableitungsbaum durchsucht sie so als \lsg[4.2cm]{Tiefensuche}.
\end{lueckentext}
\end{aufgabe}

\newpage
\begin{aufgabe}[Die Reihenfolge zählt (Website: Landkarte)]{4}
\begin{lueckentext}
Erst alle Farben, dann prüfen: \lsg[2.25cm]{7295} Schritte. Jede Prüfung so früh wie möglich: \lsg[2.25cm]{230} Schritte. Die Antworten \lsgrest{bleiben gleich (6 Färbungen).}
\end{lueckentext}
Erkläre, warum die frühe Prüfung so viel Arbeit spart.
\lsglinien{3}{Ist schon eine falsche Kombination gewählt (z. B. OB gleich NB), merkt die frühe Fassung das sofort und geht zurück. Die späte Fassung probiert vorher noch alle Farben der übrigen Bezirke durch (3 hoch 5 = 243 Möglichkeiten), die alle scheitern.}
\end{aufgabe}

\begin{aufgabe}[Endlose Rekursion (Website)]{5}
\begin{prolog}
vorfahre(X, Y) :- vorfahre(Z, Y), elternteil(X, Z).
\end{prolog}
Diese Regel ist logisch richtig. Warum wird die Inferenzmaschine trotzdem nicht fertig? Wie schreibt man rekursive Regeln richtig?
\lsglinien{4}{Sobald der Abbruchfall nichts mehr liefert, ruft die Regel als erstes Ziel wieder vorfahre(Z, Y) auf, ohne vorher einen Schritt zu gehen: Das Problem wird nie kleiner, die Suche endet nicht. Richtig: Abbruchfall zuerst, und vor dem rekursiven Aufruf ein Ziel, das das Problem verkleinert (erst elternteil, dann vorfahre), wie bei rekursiven Methoden in Java.}
\end{aufgabe}

\begin{aufgabe}[Grenzen wissensbasierter Systeme]{6}
Nenne eine Aufgabe, für die sich eine Wissensbasis gut eignet, und eine, für die nicht. Begründe.
\lsglinien{4}{Gut: Stammbaum, Fahrplan, Steuerregeln – das System ist abgeschlossen, das Wissen klar formulierbar, jede Antwort ist erklärbar. Schlecht: eine handgeschriebene 7 auf einem Bild erkennen – niemand kann die Regeln vollständig aufschreiben, Wahrnehmung ist unscharf. Dafür lernt man aus Beispielen (maschinelles Lernen).}
\end{aufgabe}
\end{document}
